\documentclass[12pt]{article}
\usepackage[utf8]{inputenc}
\usepackage[T1]{fontenc}
\usepackage{hyperref}
\title{Sarua Language Resources for K-12 Curriculum Development}
\author{Compiled by Claude Sonnet 5.5 (claude-sonnet-5-5) for Andr\'as Kornai}
\date{}
\begin{document}
\maketitle

\section{Introduction}
Sarua (also Sarwa, Saroua; ISO 639-3 swy) is an East Chadic (Afro-Asiatic) language of Chad. The only speaker figure found, about 2,000, is a 1997 SIL count and is probably out of date. Little is available for it.

\section{Primary Educational Resources}
No primers, textbooks or school materials were found.

\section{Linguistic and Reference Materials}
\begin{itemize}
\item Wikipedia overview: \url{https://en.wikipedia.org/wiki/Sarua_language}
\item Jungraithmayr and Peust (2021), a sketch of Sarwa based on original fieldwork, in \emph{Zeitschrift der Deutschen Morgenl\"andischen Gesellschaft}: \url{https://zdmg.harrassowitz-library.com/article/ZDMG/2021/1/4}
\end{itemize}

\section{Religious and Community Resources}
\begin{itemize}
\item Joshua Project profile of the Sarwa in Chad. It notes that speakers mostly use Fulfulde with outsiders and that its location differs from other sources: \url{https://joshuaproject.net/people_groups/14775}
\end{itemize}

\section{Digital and Audio Resources}
\begin{itemize}
\item Endangered Languages Project entry: \url{https://www.endangeredlanguages.com/lang/2593?hl=de}
\end{itemize}
No audio or video resources were found.

\section{Curriculum Development Context}
The language is not known to be taught in schools, and direct evidence of how it is used by children is lacking. A curriculum would need new fieldwork with the community, starting from the 2021 sketch, and would have to take into account the role of Fulfulde and other regional languages.
\end{document}